\chapter{The Entry Mode}\index{modes!Entry Mode|textbf}
\label{ch:entry}

\begin{versebox}
Five doors into every teaching ---\\
dependent arising is the key.\\
Enter through whichever gate you choose:\\
the building is the same, you see.
\end{versebox}

\noindent Think of a cathedral with five entrances. The west door, the north door, the south door, the cloister door, the crypt door. Each entrance brings you through a different corridor, past different chapels, under different arches. But every corridor leads to the same nave.\footnote{Pali: \textit{``Tattha katamo otaraṇo hāro? `Yo ca paṭiccuppādo'\,''ti.} The mnemonic fragment is ``and dependent arising'' (\textit{yo ca paṭiccuppādo}). The word \textit{otaraṇa} means ``entering, descending into, plunging into'' --- the mode that shows you how to \textit{enter} any teaching through five analytical doorways\index{entry!five doorways|textbf}: (1) the faculties (\textit{indriya}), (2) dependent origination (\textit{paṭiccasamuppāda}), (3) the aggregates (\textit{khandha}), (4) the elements (\textit{dhātu}), and (5) the sense-bases (\textit{āyatana}). Whichever door you use, you arrive at the same understanding. \textbf{A note on our approach}: the original Pali demonstrates the Entry Mode by running each passage through all five doorways in sequence, using a highly formulaic structure. We have preserved the logic and all five entries for each passage, but woven the analysis into narrative prose rather than reproducing the repetitive formula. The full Pali entries are in the endnotes.}

The Entry Mode gives you five analytical doorways into any passage the Buddha spoke. You can enter through the faculties --- the driving forces of the mind. Through dependent origination --- the chain of conditions. Through the aggregates --- the components of experience. Through the elements --- the raw building blocks of reality. Or through the sense-bases --- the meeting points of inner and outer.

Five entries. One teaching. The Entry Mode proves they all converge.

\section*{Freed in Every Direction}

The demonstration begins with a powerful verse about liberation:

\begin{versebox}
Freed above, freed below, freed everywhere ---\\
no longer seeing ``I am this'' ---\\
thus freed, he crossed the flood\\
never crossed before, to no more becoming.
\end{versebox}

``Above'' means the form and formless realms --- the higher planes of existence that meditators can access. ``Below'' means the sense-desire realm --- this world, with its sights and sounds and hungers. ``Freed everywhere'' means liberation across all three realms of existence.\footnote{Pali: \textit{``Uddhaṃ adho sabbadhi vippamutto, ayaṃ ahasmīti anānupassī; / Evaṃ vimutto udatāri oghaṃ, atiṇṇapubbaṃ apunabbhavāyā\,''ti.} (Ud 61.) \textit{``Uddha\,''nti rūpadhātu ca arūpadhātu ca. `Adho\,''ti kāmadhātu. `Sabbadhi vippamutto\,''ti tedhātuke ayaṃ asekkhāvimutti.''} ``Above'' = the form-realm (\textit{rūpadhātu}) and formless-realm (\textit{arūpadhātu}). ``Below'' = the sense-desire realm (\textit{kāmadhātu}). ``Freed everywhere in the three realms'' = the non-trainee's liberation (\textit{asekkhā vimutti}) --- the full, irreversible freedom of the arahant.}

Now watch how the Entry Mode opens its five doors.

\textit{Door one: the faculties.} The freedom described in the verse consists of the five faculties at their highest development --- faith, energy, mindfulness, concentration, and wisdom, fully matured and beyond training. These are what ``freed everywhere'' looks like from the inside: five capacities, perfected.\footnote{Pali: \textit{``Tāniyeva asekkhāni pañcindriyāni, ayaṃ indriyehi otaraṇā.''} Entry through the faculties (\textit{indriyehi otaraṇā}): ``freed everywhere'' = the five non-trainee faculties (\textit{asekkhāni pañcindriyāni}) --- faith, energy, mindfulness, concentration, wisdom at their culmination.}

\textit{Door two: dependent origination.} Those same five perfected faculties \textit{are} true knowledge. And when true knowledge arises, ignorance ceases. When ignorance ceases, formations cease. When formations cease, consciousness ceases. When consciousness ceases, name-and-form ceases. When name-and-form ceases, the six sense-bases cease. Contact ceases. Feeling ceases. Craving ceases. Clinging ceases. Existence ceases. Birth ceases. Aging-and-death, sorrow, lamentation, pain, grief, and despair --- all cease. Thus ceases this entire mass of suffering.\footnote{Pali: \textit{``Tāniyeva asekkhāni pañcindriyāni vijjā, vijjuppādā avijjānirodho, avijjānirodhā saṅkhāranirodho \ldots Evametassa kevalassa dukkhakkhandhassa nirodho hoti. Ayaṃ paṭiccasamuppādehi otaraṇā.''} Entry through dependent origination (\textit{paṭiccasamuppādehi otaraṇā}): the five faculties = true knowledge (\textit{vijjā}); the arising of true knowledge triggers the entire cessation sequence of the twelve links of dependent origination, from ignorance-cessation through to the cessation of the entire mass of suffering. One verse about freedom, read through the lens of the entire causal chain.}

The entire twelve-link chain, run in reverse. From a four-line verse about being freed.

\textit{Door three: the aggregates.} Those same five faculties are contained in three training-aggregates: virtue, concentration, and wisdom. Three baskets that hold the entire practice.\footnote{Pali: \textit{``Tāniyeva asekkhāni pañcindriyāni tīhi khandhehi saṅgahitāni --- sīlakkhandhena samādhikkhandhena paññākkhandhena, ayaṃ khandhehi otaraṇā.''} Entry through the aggregates (\textit{khandhehi otaraṇā}): the five faculties are grouped into the three training-aggregates: the aggregate of virtue (\textit{sīlakkhandha}), the aggregate of concentration (\textit{samādhikkhandha}), and the aggregate of wisdom (\textit{paññākkhandha}). Note: these are the three \textit{training}-aggregates (\textit{sikkhā}), not the five aggregates of clinging (\textit{upādānakkhandha}).}

\textit{Door four: the elements.} Those same faculties belong to the category of formations --- specifically, the taint-free formations that are not factors of continued existence. These are included in the mind-element.\footnote{Pali: \textit{``Tāniyeva asekkhāni pañcindriyāni saṅkhārapariyāpannāni ye saṅkhārā anāsavā, no ca bhavaṅgā, te saṅkhārā dhammadhātusaṅgahitā. Ayaṃ dhātūhi otaraṇā.''} Entry through the elements (\textit{dhātūhi otaraṇā}): the five faculties are formations (\textit{saṅkhāra}), specifically taint-free (\textit{anāsava}) formations that do not serve as factors of becoming (\textit{no ca bhavaṅga}). These formations are classified under the mind-object element (\textit{dhammadhātu}).}

\textit{Door five: the sense-bases.} That mind-element belongs to the mind-object sense-base --- the taint-free, non-existence-producing sense-base.\footnote{Pali: \textit{``Sā dhammadhātu dhammāyatanapariyāpannā, yaṃ āyatanaṃ anāsavaṃ, no ca bhavaṅgaṃ. Ayaṃ āyatanehi otaraṇā.''} Entry through the sense-bases (\textit{āyatanehi otaraṇā}): the mind-object element (\textit{dhammadhātu}) is included in the mind-object sense-base (\textit{dhammāyatana}), specifically the taint-free, non-becoming-factor portion. This completes the five-door analysis.}

Five doors. Five analytical frameworks. Same verse, same freedom, same liberation --- entered from five different angles.

\section*{The End of ``I Am''}

But the verse has a second half: ``No longer seeing `I am this.'\,'' That phrase gets its own five-door analysis, because it describes a different aspect of liberation --- the uprooting of identity-view, which is the trainee's liberation, the breakthrough at stream-entry.\footnote{Pali: \textit{``Ayaṃ ahasmīti anānupassī\,''ti ayaṃ sakkāyadiṭṭhiyā samugghāto, sā sekkhāvimutti, tāniyeva sekkhāni pañcindriyāni. Ayaṃ indriyehi otaraṇā.''} ``No longer seeing `I am this'\,'' = the uprooting of identity-view (\textit{sakkāyadiṭṭhi}), which is the \textit{trainee's} liberation (\textit{sekkhā vimutti}), as distinct from the non-trainee's full liberation discussed above. The Netti runs this phrase through all five doorways again: (1) the five \textit{trainee} faculties; (2) true knowledge → cessation of ignorance → cessation of the entire chain of dependent origination; (3) the wisdom aggregate; (4) taint-free formations in the mind-object element; (5) taint-free mind-object sense-base.}

The faculties here are the \textit{trainee's} five faculties --- not yet perfected, but awakened. The same dependent origination chain runs in reverse. The same aggregates, elements, and sense-bases apply. But the level is different. The trainee has crossed the threshold; the non-trainee has completed the journey.

And the verse's final line --- ``thus freed, he crossed the flood never crossed before, to no more becoming'' --- combines both liberations, trainee and non-trainee, into the single accomplishment of one who has finished the work.\footnote{Pali: \textit{``Sekkhāya ca vimuttiyā asekkhāya ca vimuttiyā vimutto udatāri oghaṃ atiṇṇapubbaṃ apunabbhavāya.''} The crossing of the flood encompasses both liberations: the trainee's breakthrough (\textit{sekkhā vimutti}) and the non-trainee's completion (\textit{asekkhā vimutti}). Together, they form the complete passage from bondage to freedom --- ``to no more becoming'' (\textit{apunabbhavāya}).}

\section*{The Dependent One Is Shaken}

The second demonstration uses one of the most compressed passages in the Canon:

\begin{quote}
The dependent one is shaken. The non-dependent one is not shaken. When there is no shaking, there is tranquility. When there is tranquility, there is no inclination. When there is no inclination, there is no coming-and-going. When there is no coming-and-going, there is no passing-away-and-reappearing. When there is no passing-away-and-reappearing --- neither here, nor there, nor in between. Just this is the end of suffering.
\end{quote}

A chain of cessation in eight links. Each link removes one more layer of bondage until nothing remains but the end of suffering.\footnote{Pali: \textit{``Nissitassa calitaṃ, anissitassa calitaṃ natthi, calite asati passaddhi, passaddhiyā sati nati na hoti, natiyā asati āgatigati na hoti, āgatigatiyā asati cutūpapāto na hoti, cutūpapāte asati nevidha na huraṃ na ubhayamantarena esevanto dukkhassā\,''ti.} (Ud 74.) Eight stages of cessation: dependence → shaking → [no tranquility] → inclination → coming-and-going → passing-away-and-reappearing → being ``here'' or ``there'' or ``in between'' → suffering. Remove the first, and all the rest collapse like dominoes.}

The Entry Mode begins with the first link. ``The dependent one is shaken'' --- but what is ``dependence''? Two kinds: dependence on craving and dependence on views. The lustful person's volition is craving-dependence. The deluded person's volition is view-dependence.\footnote{Pali: \textit{``Nissayo nāma duvidho taṇhānissayo ca diṭṭhinissayo ca. Tattha yā rattassa cetanā, ayaṃ taṇhānissayo; yā mūḷhassa cetanā, ayaṃ diṭṭhinissayo.''} Two types of dependence (\textit{nissaya}): craving-dependence (\textit{taṇhā-nissaya}) and view-dependence (\textit{diṭṭhi-nissaya}). The lustful person's (\textit{ratta}) volition (\textit{cetanā}) = craving-dependence; the confused person's (\textit{mūḷha}) volition = view-dependence. Both are forms of bondage; both produce shaking.}

Now the five doors open. Volition is a formation. Formations condition consciousness. Consciousness conditions name-and-form. And the whole chain of dependent origination unfolds. That is door two.\footnote{Pali: \textit{``Cetanā pana saṅkhārā, saṅkhārapaccayā viññāṇaṃ, viññāṇapaccayā nāmarūpaṃ, evaṃ sabbo paṭiccasamuppādo. Ayaṃ paṭiccasamuppādehi otaraṇā.''} Entry through dependent origination: volition (\textit{cetanā}) = formations (\textit{saṅkhārā}); formations → consciousness → name-and-form → and so on through the full twelve-link chain.}

The lustful person's feeling is pleasant; the deluded person's feeling is neutral. Two feelings --- and feeling is an aggregate. That is door three.\footnote{Pali: \textit{``Tattha yā rattassa vedanā, ayaṃ sukhā vedanā. Yā sammūḷhassa vedanā, ayaṃ adukkhamasukhā vedanā, imā dve vedanā vedanākkhandho. Ayaṃ khandhehi otaraṇā.''} Entry through the aggregates: the lustful person's feeling = pleasant feeling; the deluded person's feeling = neither-painful-nor-pleasant feeling. Both belong to the feeling aggregate (\textit{vedanākkhandha}).}

Pleasant feeling corresponds to the pleasure-faculty and the joy-faculty; neutral feeling corresponds to the equanimity-faculty. Three faculties. That is door one.\footnote{Pali: \textit{``Tattha sukhā vedanā dve indriyāni sukhindriyaṃ somanassindriyañca, adukkhamasukhā vedanā upekkhindriyaṃ. Ayaṃ indriyehi otaraṇā.''} Entry through the faculties: pleasant feeling = the pleasure-faculty (\textit{sukhindriya}) and the joy-faculty (\textit{somanassindriya}); neutral feeling = the equanimity-faculty (\textit{upekkhindriya}).}

These faculties are formations --- taint-laden formations that serve as factors of continued existence, classified under the mind-element. Door four. And that mind-element belongs to the taint-laden mind-object sense-base. Door five.\footnote{Pali: \textit{``Tāniyeva indriyāni saṅkhārapariyāpannāni, ye saṅkhārā sāsavā bhavaṅgā, te saṅkhārā dhammadhātusaṅgahitā. Ayaṃ dhātūhi otaraṇā. Sā dhammadhātu dhammāyatanapariyāpannā, yaṃ āyatanaṃ sāsavaṃ bhavaṅgaṃ. Ayaṃ āyatanehi otaraṇā.''} Entry through elements and sense-bases. Note the contrast with the ``freed everywhere'' verse: there, the formations were taint-\textit{free} and \textit{not} factors of becoming; here, they are taint-\textit{laden} and \textit{are} factors of becoming (\textit{sāsavā bhavaṅgā}). The Entry Mode is sensitive to this distinction --- the same five-door framework produces different results depending on whether the passage describes bondage or freedom.}

Notice the crucial difference from the previous verse. There, the formations were taint-free. Here, they are taint-laden. The five-door framework is identical, but the content --- bondage versus freedom --- reverses the polarity entirely.

\section*{``Not Shaken''}

``The non-dependent one is not shaken.'' How does one become non-dependent? Through calm, one is freed from dependence on craving\index{liberation}. Through insight, one is freed from dependence on views\index{view-character}.\footnote{Pali: \textit{``Anissitassa calitaṃ natthī\,''ti samathavasena vā taṇhāya anissito vipassanāvasena vā diṭṭhiyā anissito.''} Two antidotes for two types of dependence: calm (\textit{samatha}) frees from craving-dependence; insight (\textit{vipassanā}) frees from view-dependence.}

And insight \textit{is} true knowledge. The five doors open again, this time in reverse: true knowledge arises, ignorance ceases, the whole chain collapses. Insight is the wisdom-aggregate --- door three. Insight corresponds to the energy-faculty and the wisdom-faculty --- door one. Insight is a taint-free formation in the mind-element --- door four. That element belongs to the taint-free mind-object sense-base --- door five.\footnote{Pali: \textit{``Yā vipassanā ayaṃ vijjā, vijjuppādā avijjānirodho \ldots Ayaṃ paṭiccasamuppādehi otaraṇā. Sāyeva vipassanā paññākkhandho. Ayaṃ khandhehi otaraṇā. Sāyeva vipassanā dve indriyāni --- vīriyindriyañca paññindriyañca. Ayaṃ indriyehi otaraṇā. Sāyeva vipassanā saṅkhārapariyāpannā, ye saṅkhārā anāsavā, no ca bhavaṅgā \ldots Ayaṃ dhātūhi otaraṇā \ldots Ayaṃ āyatanehi otaraṇā.''} The five-door analysis for ``not shaken'': (1) dependent origination: insight = true knowledge → ignorance ceases → the entire chain ceases; (2) aggregates: insight = the wisdom aggregate; (3) faculties: insight = energy-faculty + wisdom-faculty; (4) elements: insight = taint-free formations in the mind-object element; (5) sense-bases: the mind-object element = taint-free mind-object sense-base. Same five doors, now showing the path of liberation.}

The same building, entered from the other side.

\section*{Tranquility, All the Way Down}

``When there is tranquility \ldots'' The passage now describes what happens after non-dependence is established. Tranquility is twofold: bodily and mental. And from tranquility comes the most remarkable chain in all of Buddhist psychology:

A tranquil body experiences pleasure. A person in pleasure becomes concentrated. A concentrated person understands things as they really are. Understanding, one becomes disenchanted. Disenchanted, one fades away from attachment. Through fading away, one is liberated. In liberation comes the knowledge: ``Birth is destroyed. The holy life has been lived. What had to be done has been done. There is no more of this.''\footnote{Pali: \textit{``Passaddhakāyo sukhaṃ vediyati, sukhino cittaṃ samādhiyati, samāhito yathābhūtaṃ pajānāti, yathābhūtaṃ pajānanto nibbindati, nibbindanto virajjati, virāgā vimuccati, vimuttasmiṃ `vimutta'miti ñāṇaṃ hoti, `khīṇā jāti, vusitaṃ brahmacariyaṃ, kataṃ karaṇīyaṃ, nāparaṃ itthattāyā'\,''ti pajānāti.''} The famous sequence: tranquility → pleasure → concentration → seeing-things-as-they-are → disenchantment → fading of passion → liberation → reviewing knowledge. This is one of the most frequently recurring formulas in the Canon (cf.\ AN 10.2, AN 11.2). Each step follows naturally from the one before --- not by force, but by an inner logic: the body calms, pleasure arises, concentration follows pleasure, clarity follows concentration, disenchantment follows clarity, release follows disenchantment.}

And then the liberated person no longer inclines toward forms, sounds, smells, tastes, touches, or ideas --- because lust, hatred, and delusion have been destroyed. And here the Netti quotes one of the most philosophically daring passages in the Canon:

\begin{quote}
Whatever form by which one might describe the Tathāgata, standing or walking --- that form has been destroyed, faded, ceased, relinquished, let go. The Tathāgata cannot be said to ``exist,'' cannot be said to ``not exist,'' cannot be said to ``both exist and not exist,'' cannot be said to ``neither exist nor not exist.'' Rather: deep, immeasurable, unfathomable --- quenched.
\end{quote}

And the same for feeling, perception, formations, and consciousness. Whatever you might use to describe the fully awakened one --- it has been abandoned. The awakened one cannot be captured in any category of existence or non-existence. All that remains is the word ``quenched'' --- and even that is just a conventional designation for what lies beyond all designation.\footnote{Pali: \textit{``So na namati rūpesu, na saddesu \ldots yena rūpena tathāgataṃ tiṭṭhantaṃ carantaṃ paññāpayamāno paññāpeyya, tassa rūpassa khayā virāgā nirodhā cāgā paṭinissaggā rūpasaṅkhaye vimutto, tathāgato atthītipi na upeti, natthītipi na upeti, atthi natthītipi na upeti, nevatthi no natthītipi na upeti. Atha kho gambhīro appameyyo asaṅkheyyo nibbutotiyeva saṅkhaṃ gacchati khayā rāgassa, khayā dosassa, khayā mohassa.''} This passage recurs throughout the Canon (cf.\ SN 44.1, MN 72). The four propositions that the Tathāgata transcends --- existence, non-existence, both, neither --- constitute the classical \textit{catuṣkoṭi} (``four corners'') of Indian logic. The Tathāgata cannot be located in any of them. The term \textit{nibbutotiyeva saṅkhaṃ gacchati} --- ``is reckoned merely as `quenched'\,'' --- suggests that even the designation ``quenched'' (\textit{nibbuta}) is a concession to conventional language, not a description of an ultimate reality. The Netti presents this passage as the final horizon of the Entry Mode: the point where all five analytical doorways reach a reality that can no longer be entered by analysis at all.}

\section*{The Source of Sorrow}

The Netti applies the Entry Mode to two more passages. First, a verse about love and loss:

\begin{versebox}
Whatever sorrows, whatever lamentations,\\
whatever sufferings of many kinds exist in the world ---\\
all arise dependent on what is dear.\\
When there is nothing dear, these do not exist.

\stanza

Therefore happy and free from sorrow\\
are those for whom nothing in the world is dear.\\
Therefore, aspiring to the sorrowless and stainless,\\
make nothing dear to you anywhere in the world.
\end{versebox}

The first stanza describes the problem; the second describes the cure. And the Entry Mode enters each through all five doors.\footnote{Pali: \textit{``Ye keci sokā paridevitā vā, dukkhā ca lokasmimanekarūpā; / Piyaṃ paṭiccappabhavanti ete, piye asante na bhavanti ete. / Tasmā hi te sukhino vītasokā, yesaṃ piyaṃ natthi kuhiñci loke; / Tasmā asokaṃ virajaṃ patthayāno, piyaṃ na kayirātha kuhiñci loke\,''ti.} (Ud 78.) The Netti splits the verse in two: the first stanza (the problem) is analyzed as containing painful feeling (\textit{dukkhā vedanā}) and pleasant feeling (\textit{sukhā vedanā}) --- ``sorrows arise dependent on what is dear'' involves suffering; ``when there is nothing dear, these do not exist'' involves pleasure (freedom from sorrow). Both feelings belong to the feeling aggregate → entry through aggregates. Feeling conditions craving, craving conditions clinging, and so on → entry through dependent origination. Pleasant feeling = pleasure-faculty + joy-faculty; painful feeling = pain-faculty + displeasure-faculty → entry through faculties. These are taint-laden formations in the mind-object element → entry through elements. The mind-object element belongs to the taint-laden sense-base → entry through sense-bases.}

``All arise dependent on what is dear'' --- the sorrow side. Sorrow is painful feeling; the relief when what is dear is absent is pleasant feeling. Both are part of the feeling-aggregate. That is the aggregates door. Feeling conditions craving, craving conditions clinging, clinging conditions existence --- the full dependent origination chain. That is the dependent origination door. Pleasant feeling maps to the pleasure- and joy-faculties; painful feeling maps to the pain- and displeasure-faculties. The faculties door. And so on through elements and sense-bases.

The second stanza --- ``aspiring to the sorrowless and stainless, make nothing dear'' --- is the abandonment of craving. The Netti identifies this as calm: the practice of releasing attachment. And calm is the concentration-aggregate, the mindfulness-faculty and concentration-faculty, taint-free formations in the mind-element, the taint-free mind-object sense-base.\footnote{Pali: \textit{``Idaṃ taṇhāpahānaṃ. Taṇhānirodhā upādānanirodho \ldots Ayaṃ paṭiccasamuppādehi otaraṇā. Taṃyeva taṇhāpahānaṃ samatho. So samatho dve indriyāni satindriyaṃ samādhindriyañca. Ayaṃ indriyehi otaraṇā. Soyeva samatho samādhikkhandho. Ayaṃ khandhehi otaraṇā \ldots''} The second stanza's analysis: abandoning craving → craving ceases → clinging ceases → the entire chain ceases (dependent origination door). This abandonment = calm (\textit{samatha}). Calm = mindfulness-faculty + concentration-faculty (faculties door). Calm = concentration-aggregate (aggregates door). Calm = taint-free formations in the mind-object element (elements door). The mind-object element = taint-free mind-object sense-base (sense-bases door).}

Five doors into sorrow. Five doors out.

\section*{The Desire That Destroys}

One final demonstration. A verse about the trap of getting what you want:

\begin{versebox}
When a person desires a pleasure\\
and that desire is fulfilled ---\\
certainly that mortal is delighted,\\
having gotten what he wished for.

\stanza

But when those pleasures fade\\
for that person full of desire and longing ---\\
he is crushed,\\
as if pierced by an arrow.
\end{versebox}

Delight when you get what you want --- that is attraction, the pull toward pleasure. Anguish when you lose it --- that is aversion, the recoil from pain. Both are sides of the same coin, and both belong to the craving-side of the mind.\footnote{Pali: \textit{``Kāmaṃ kāmayamānassa, tassa ce taṃ samijjhati; / Addhā pītimano hoti, laddhā macco yadicchati. / Tassa ce kāmayānassa, chandajātassa jantuno; / Te kāmā parihāyanti, sallaviddhova ruppati.''} (Sn 766--767.) \textit{``Tattha yā pītimanatā, ayaṃ anunayo. Yadāha sallaviddhova ruppatīti, idaṃ paṭighaṃ. Anunayaṃ paṭighañca pana taṇhāpakkho.''} The Netti identifies two mental reactions: attraction (\textit{anunaya}) when desire is fulfilled, and aversion (\textit{paṭigha}) when pleasures fade. Both belong to the craving-side (\textit{taṇhāpakkha}) of the mind. Both are fuel for continued existence.}

And the basis of craving? The ten physical sense-bases --- the five sense-organs and their five objects. That is the sense-bases door. Those ten physical bases, combined with the mental body, form name-and-form --- and from name-and-form, the six sense-bases arise, then contact, then feeling, then craving, and the whole chain rolls on. That is the dependent origination door. Name-and-form is the five aggregates. The aggregates door. Name-and-form is the eighteen elements. The elements door. The physical body maps to five physical faculties; the mental body maps to five non-physical faculties --- ten in total. The faculties door.\footnote{Pali: \textit{``Taṇhāya ca pana dasarūpīni āyatanāni padaṭṭhānaṃ. Ayaṃ āyatanehi otaraṇā. Tāniyeva dasa rūpīni rūpakāyo nāmasampayutto, tadubhayaṃ nāmarūpaṃ \ldots Ayaṃ paṭiccasamuppādehi otaraṇā. Tadeva nāmarūpaṃ pañcakkhandho. Ayaṃ khandhehi otaraṇā. Tadeva nāmarūpaṃ aṭṭhārasa dhātuyo. Ayaṃ dhātūhi otaraṇā. Tattha yo rūpakāyo imāni pañca rūpīni indriyāni, yo nāmakāyo imāni pañca arūpīni indriyāni, imāni dasa indriyāni. Ayaṃ indriyehi otaraṇā.''} The five-door entry for the desire verse: (1) sense-bases: craving's basis = the ten physical sense-bases; (2) dependent origination: name-and-form → six sense-bases → contact → feeling → craving → and so on; (3) aggregates: name-and-form = the five aggregates; (4) elements: name-and-form = the eighteen elements; (5) faculties: the physical body = five physical faculties, the mental body = five non-physical faculties.}

And the cure? The verse continues:

\begin{versebox}
Whoever avoids pleasures\\
as one would avoid stepping on a serpent's head ---\\
that mindful one transcends\\
this clinging in the world.
\end{versebox}

This is the nirvana-element with residue remaining --- the living liberation of one who still walks the earth but no longer clings. And this too gets its five-door analysis: it \textit{is} true knowledge (dependent origination door), it \textit{is} the wisdom-aggregate (aggregates door), it maps to the energy- and wisdom-faculties (faculties door), it is taint-free formations in the mind-element (elements door), and that element belongs to the taint-free mind-object sense-base (sense-bases door).\footnote{Pali: \textit{``Yo kāme parivajjeti, sappasseva padā siro; / Somaṃ visattikaṃ loke, sato samativattatī\,''ti. Ayaṃ saupādisesā nibbānadhātu \ldots Sāyeva saupādisesā nibbānadhātu vijjā, vijjuppādā avijjānirodho \ldots Sāyeva vijjā paññākkhandho \ldots Sāyeva vijjā dve indriyāni --- vīriyindriyaṃ paññindriyañca \ldots Sāyeva vijjā saṅkhārapariyāpannā, ye saṅkhārā anāsavā, no ca bhavaṅgā \ldots Sā dhammadhātu dhammāyatanapariyāpannā, yaṃ āyatanaṃ anāsavaṃ, no ca bhavaṅgaṃ. Ayaṃ āyatanehi otaraṇā.''} The nirvana-element with residue remaining (\textit{saupādisesā nibbānadhātu}) = true knowledge → dependent origination door. True knowledge = wisdom aggregate → aggregates door. True knowledge = energy-faculty + wisdom-faculty → faculties door. True knowledge = taint-free formations in the mind-object element → elements door. The mind-object element = taint-free mind-object sense-base → sense-bases door.}

\ornament

The Entry Mode's message is simple and radical. Every teaching the Buddha gave --- no matter how poetic, how compressed, how seemingly straightforward --- can be entered through all five of these analytical doorways. The faculties show you the inner capacities involved. Dependent origination shows you the causal chain. The aggregates show you the components of experience. The elements show you the raw material. The sense-bases show you where inner meets outer.

And every doorway leads to the same nave: the understanding that whatever is dependently arisen is subject to cessation, and that cessation is the end of suffering.\footnote{Pali: \textit{``Ettāvatā paṭicca indriyakhandhadhātuāyatanāni samosaraṇotaraṇāni bhavanti. Evaṃ paṭicca indriyakhandhadhātuāyatanāni otāretabbāni. Tenāha āyasmā mahākaccāyano `yo ca paṭiccuppādo'\,''ti.''} The closing formula: ``To this extent, the faculties, aggregates, elements, and sense-bases converge as entries. Thus should [passages] be entered through the faculties, aggregates, elements, and sense-bases. Therefore the Venerable Mahākaccāyana said: `and dependent arising.'\,'' The word \textit{samosaraṇa} (``converging'') is key: five different analytical pathways, one point of convergence.}
